% !TeX root = ../../packet-and-principal-models.tex
\section{The explicit polarization on the normalized Hilbert modular model}\label{sec:polarization}
The preceding arguments use $X_L=\mathbb H^6/\mathrm{SL}_2(\OL)$ as an analytic Hilbert modular quotient. To apply a theorem about polarized abelian varieties we specify the polarization directly rather than infer principal polarization from the notation $X_L$.

For $z=(z_1,\ldots,z_6)\in\mathbb H^6$, let $\sigma_1,\ldots,\sigma_6$ be the real embeddings of $L$ and set
\[
 \Lambda_z=\{(\sigma_i(a)z_i+\sigma_i(b))_{i=1}^6:a,b\in\OL\}\subset\mathbb C^6,
 \qquad A_z=\mathbb C^6/\Lambda_z.
\]
On the abstract lattice $\OL^2$ define
\begin{equation}\label{eq:E1}
 E_1((a,b),(a',b'))=\operatorname{Tr}_{L/\Ql}(a'b-ab').
\end{equation}
This form is integral and alternating. With the displayed sign convention, at one real place, if $v=(a,b)$ and $J_z$ is multiplication by $i$ in the coordinate $az+b$, then
\[
 E_1(v,J_zv)=a^2\,\operatorname{Im}(z)+\frac{(a\operatorname{Re}(z)+b)^2}{\operatorname{Im}(z)}>0
\]
for $v\neq0$. Summing over the six embeddings verifies the Riemann bilinear positivity criterion, so $E_1$ is a polarization. The $\OL$-action is self-adjoint for $E_1$, and every matrix in $\mathrm{SL}_2(\OL)$ preserves \eqref{eq:E1}. Thus this construction descends to the quotient stack $[\mathbb H^6/\mathrm{SL}_2(\OL)]$ (equivalently, to the Hilbert modular orbifold $X_L$) as a polarized RM family.  After passage to any neat finite-index subgroup it becomes an honest abelian scheme over the corresponding finite cover.  In the arguments below only the explicit fiber at $z_B$ is needed.

\begin{proposition}[Polarization type]\label{prop:E1type}
For the integral basis $1,t,\ldots,t^5$ of $\OL$, the polarization $E_1$ has type
\[
 \boxed{(1,7,7,21,21,21)}.
\]
In particular, the polarization used below is not principal. Its degree is
\[
 (1\cdot7\cdot7\cdot21\cdot21\cdot21)^2=\operatorname{disc}(L)^2.
\]
\end{proposition}
\begin{proof}
Let
\[
 T=(\operatorname{Tr}_{L/\Ql}(t^{i+j}))_{0\le i,j\le5}.
\]
In the induced basis of $\OL^2$, the alternating matrix of $E_1$ is
$\left(\begin{smallmatrix}0&-T\\T&0\end{smallmatrix}\right)$. Exact Smith reduction gives
\[
 \operatorname{SNF}(T)=\operatorname{diag}(1,7,7,21,21,21),
 \qquad \det T=3^3 7^5=453789=\operatorname{disc}(L).
\]
The elementary divisors of the alternating form are therefore the six displayed entries, and the polarization degree is the square of their product.
\end{proof}

\begin{proposition}[No compatible principal polarization on the explicit $\OL^2$ lattice]\label{prop:noprincipal}
For the fixed real-multiplication structure on $H_1(A_z,\Z)\cong\OL^2$, there is no principal polarization whose Rosati involution fixes $L$ pointwise.
\end{proposition}
\begin{proof}
An $L$-self-adjoint alternating form on the rank-two $L$-space $L^2$ is necessarily of the form
\[
 E_\lambda((a,b),(a',b'))=\Tr_{L/\Ql}\!\left(\lambda(a'b-ab')\right)
\]
for a unique $\lambda\in L$.  On the lattice $\OL^2$ it is integral exactly when $\lambda\in\mathfrak D_L^{-1}$, and it is unimodular exactly when $\lambda\mathfrak D_L=\OL$.  With $z\in\HH^6$, the Riemann positivity condition is $\sigma_i(\lambda)>0$ for all six real embeddings.  Thus a compatible principal polarization exists precisely when the codifferent has a totally positive generator.

Because $\OL=\Z[t]$ with minimal polynomial $f$, the different is $(f'(t))$.  In the embedding order $k=1,5,11,13,17,19$, the signs of $f'(\sigma_k(t))$ are
\[
 (+,-,+,-,+,-).
\]
Supplement, Section~\externalref{supplement}{app:orders} computes the image of the unit-sign map $E_L\to\{\pm1\}^6$: it is the rank-five subgroup generated by the sign vectors of $\xi_5,\xi_{11},\xi_{13},\xi_{17},\xi_{19}$ and $-1$.  Exact row reduction over $\mathbf F_2$ shows that $(+,-,+,-,+,-)$ is not in this subgroup.  Hence no unit multiple of $1/f'(t)$ is totally positive, so $\mathfrak D_L^{-1}$ has no totally positive generator.
\end{proof}

\paragraph{Homology convention.}
The column $(a,b)$ represents $az+b$.  If a Mobius matrix
$M=\left(\begin{smallmatrix}r&s\\u&v\end{smallmatrix}\right)$ fixes $z$,
then multiplication by $uz+v$ on these columns is $M^{\mathsf T}$, as
proved in Proposition~\ref{prop:fixedmatrixlattice}.  Accordingly, a
Mobius-polynomial representative $S_X$ of $\sqrt{-3}$ acts on homology as
$\Sigma_X=S_X^{\mathsf T}$.  All lattice stability and Hermitian forms
below use $\Sigma_X$.

\begin{proposition}[A principally polarized model in the isogeny class]\label{prop:ppisogenous}
Let $\Lambda=\mathcal O_L^2$ be the homology lattice of $A_B$, let
\[
 \Lambda^\#=\{v\in L^2:E_1(v,\Lambda)\subset\mathbb Z\}
              =\mathfrak D_L^{-1}\Lambda,
\]
and let $\Sigma_B=(I+2B_t^{14})^{\mathsf T}$.  There is an $\mathcal O_L$-submodule
\[
 H\subset \Lambda^\#/\Lambda
\]
which is maximal isotropic for the discriminant pairing and stable under the action of $\mathcal O_K$.  Here $\sqrt{-3}$ acts as $\Sigma_B$, while the integral generator
\[
 U:=(B_t^{14})^{\mathsf T}=\frac{\Sigma_B-I}{2}
\]
satisfies $U^2+U+I=0$ and hence identifies $\mathcal O_K$ with $\mathbb Z[U]$.  If $\Lambda^{\mathrm{pp}}$ denotes the inverse image of $H$ in $\Lambda^\#$, then
\[
 [\Lambda^{\mathrm{pp}}:\Lambda]=\operatorname{disc}(L)=453789,
 \qquad (\Lambda^{\mathrm{pp}})^\#=\Lambda^{\mathrm{pp}}.
\]
Consequently
\[
 A_B^{\mathrm{pp}}:=\mathbb C^6/\Lambda_B^{\mathrm{pp}}
\]
is a principally polarized abelian sixfold, the natural map
\[
 A_B\longrightarrow A_B^{\mathrm{pp}}
\]
is an $\mathcal O_L$-linear isogeny of degree $453789$, and the $K$-action descends to $A_B^{\mathrm{pp}}$.  The rational Hodge structure, $K$-signature $(3,3)$, Weil discriminant $-1$, and validity of the Hodge conjecture are unchanged under this isogeny.
\end{proposition}
\begin{proof}
Since $\mathcal O_L=\mathbb Z[t]$, Dedekind factorization applied to~\eqref{eq:minpoly} gives
\[
 f(T)\equiv (T^3-T^2+T+1)^2\pmod 3,
 \qquad
 f(T)\equiv (T-1)^6\pmod 7.
\]
Both ramification indices are prime to the residue characteristic, hence the ramification is tame.  Thus
\[
 (3)=\mathfrak p_3^2,\quad N\mathfrak p_3=3^3,
 \qquad
 (7)=\mathfrak p_7^6,\quad N\mathfrak p_7=7,
\]
and the different is
\[
 \mathfrak D_L=\mathfrak p_3\mathfrak p_7^5.
\]
Multiplication by a generator of the different identifies the discriminant module with
\[
 \Lambda^\#/\Lambda\cong
 (\mathcal O_L/\mathfrak D_L)^2
 \cong
 (\mathcal O_L/\mathfrak p_3)^2\oplus
 (\mathcal O_L/\mathfrak p_7^5)^2.
\]

At $\mathfrak p_3$, exact reduction of the displayed matrix $\Sigma_B$ gives
\[
 \Sigma_B\equiv0\pmod{\mathfrak p_3}.
\]
Hence any one-dimensional $\mathcal O_L/\mathfrak p_3\cong\mathbb F_{27}$ subspace is $\Sigma_B$-stable; take the first-coordinate line $H_3$.  Since $U=(\Sigma_B-I)/2\equiv1\pmod{\mathfrak p_3}$, this line is already $\mathcal O_K=\mathbb Z[U]$-stable.

At $\mathfrak p_7^5$, the quotient has characteristic $7$ and $\Sigma_B^2=-3I=4I$.  Moreover $U^2+U+I=0$ gives $\operatorname{tr}(U)=-1$, hence
\[
 \operatorname{tr}(\Sigma_B)=\operatorname{tr}(I+2U)=0.
\]
Since $4$ is a unit, the idempotents
\[
 e_+=\frac{\Sigma_B+2I}{4},\qquad e_-=\frac{2I-\Sigma_B}{4}
\]
split $(\mathcal O_L/\mathfrak p_7^5)^2$ into the $+2$ and $-2$ eigensummands.  Because the local ring $\mathcal O_L/\mathfrak p_7^5$ is local, the images of these idempotents are finite projective and hence free.  Moreover
\[
 \operatorname{tr}(e_+)=\frac{\operatorname{tr}(\Sigma_B)+4}{4}=1,
 \qquad
 \operatorname{tr}(e_-)=1,
\]
so both images have rank one.  Let $H_7=\operatorname{im}(e_+)$.  Because $\Sigma_B^\dagger=-\Sigma_B$, for $v,w\in H_7$ one has
\[
 0=E_1(\Sigma_Bv,w)+E_1(v,\Sigma_Bw)=4E_1(v,w)
\]
in the local discriminant pairing, so $H_7$ is isotropic.  The line $H_3$ is isotropic for the same alternating pairing.  By the Chinese remainder theorem
\[
 H=H_3\oplus H_7
\]
is isotropic and $\mathcal O_L$-stable, and
\[
 |H|=3^3\,7^5=453789=\operatorname{disc}(L)
       =\sqrt{|\Lambda^\#/\Lambda|}.
\]
Thus $H$ is maximal isotropic and its inverse image $\Lambda^{\mathrm{pp}}$ is self-dual for $E_1$.

The action of $\Sigma_B$ preserves $H$.  The full ring $\mathcal O_K=\mathbb Z[U]$ preserves it as well.  Indeed $U=(B_t^{14})^{\mathsf T}$ preserves both $\Lambda$ and $\Lambda^\#$; on the discriminant module one also has $U=(\Sigma_B-I)/2$, and multiplication by $2$ is invertible on both the $3$- and $7$-primary parts, so $\Sigma_B$-stability of $H$ implies $U$-stability.  Hence $\Lambda^{\mathrm{pp}}$ is $\mathcal O_K$-stable.  The complex structure, the rational $K$-action on $V:=\Lambda\otimes\Ql=L^2$, and the rational alternating form $E_1$ have not changed.  Hence the $K$-Hermitian form $H(v,w)=E_1(\Sigma_Bv,w)+\sqrt{-3}\,E_1(v,w)$ on the rational $K$-space $V$ is literally the same form before and after the lattice enlargement.  In particular its determinant class in $\Ql^\times/N_{K/\Ql}(K^\times)$ remains $-1$, not merely up to an appeal to isogeny invariance.  The CM type and all rational Hodge classes are likewise unchanged.  The inclusion $\Lambda\subset\Lambda^{\mathrm{pp}}$ has index $|H|$, which is the degree of the induced isogeny.  Finally the Hodge conjecture is invariant under isogeny (pullback and pushforward along a finite isogeny preserve algebraicity and induce inverse isomorphisms on rational cohomology up to the degree).
\end{proof}

\begin{corollary}[The principal model is not the same orbifold vertex]\label{cor:noBstablepp}
No principally polarized self-dual overlattice
\[
 \Lambda\subset\Lambda'\subset\Lambda^\#
\]
can be stable under the full vertex automorphism acting as $B_t^{\mathsf T}$ on homology.  In particular the principally polarized model of Proposition~\ref{prop:ppisogenous} cannot be obtained by merely changing the polarization while retaining the same integral order-$42$ $B_t$-vertex structure.
\end{corollary}
\begin{proof}
Every self-dual overlattice corresponds to a maximal isotropic subgroup,
whose $\mathfrak p_3$-part has order $27$. Stability under the actual
homology automorphism $B_t^{\mathsf T}$ implies stability under its inverse
and therefore under their sum $tI$. Hence this subgroup is an
$\mathbb F_{27}$-line in $(\mathcal O_L/\mathfrak p_3)^2$.  The characteristic polynomial of $B_t$ is $X^2-tX+1$.  Its discriminant $t^2-4$ is a nonsquare in $\mathbb F_{27}$: exact reduction modulo $g_3(T)=T^3-T^2+T+1$ gives
\[
 (t^2-4)^{13}=-1.
\]
Hence $X^2-tX+1$ is irreducible over $\mathbb F_{27}$ and $B_t$ has no invariant $\mathbb F_{27}$-line.  Thus no maximal isotropic subgroup, and hence no principally polarized overlattice, can be $B_t^{\mathsf T}$-stable.
\end{proof}

\begin{remark}[From special fibers to a common principal family]\label{rem:ppcover}
Proposition~\ref{prop:commonprincipal} below supplies a common nonfree
maximal isotropic subgroup for the three specialized $K$-actions.
Proposition~\ref{prop:principalstabilizer} determines its ambient index
$224$ and its index $28$ in the explicitly parametrizing quadrilateral
group, before auxiliary level. Theorem~\ref{thm:principalfamily} descends
its inverse-image lattice sheaf and principal polarization over a connected
neat cover. Proposition~\ref{prop:survivingorder} determines the surviving
order inside each selected CM field. The value $537824$ in the earlier
free-line construction concerns a different subgroup and is not the
index of this common nonfree level.
\end{remark}

\begin{remark}[What the principally polarized model does not prove]\label{rem:pphec}
The special-fiber subgroup $H$ gives a principally polarized model
$A_B^{\mathrm{pp}}$; Remark~\ref{rem:ppcover} describes the later common
choice used for the family. Corollary~\ref{cor:noBstablepp}
shows that any such modification necessarily loses the integral order-$42$
vertex automorphism at the prime above $3$, so it would not be the original
orbifold $B$-vertex in another polarization disguise.  We do not identify
$A_B^{\mathrm{pp}}$ with a point on Mostaed's original curve $V$ inside the
principally polarized Hilbert--Blumenthal moduli space.
\end{remark}

Thus the Hodge-theoretic conclusion below does not require a map from $X_L$ to the principally polarized moduli space $\mathcal A_6$: it uses the explicit polarized abelian variety $(A_z,E_1)$ at the certified vertex.  In fact Proposition~\ref{prop:noprincipal} shows that, for this fixed $\OL$-module and real-multiplication action, the missing principal polarization is not merely unconstructed.  Markman's theorem is stated for polarized abelian sixfolds, not only principally polarized ones.

At the normalized $B$-vertex put
\[
 U=(B_t^{14})^{\mathsf T},\qquad \Sigma_B=I+2U.
\]
Since $B_t^{\mathsf T}\in\mathrm{SL}_2(\OL)$, its $E_1$-adjoint is $(B_t^{-1})^{\mathsf T}$. The elliptic relation gives $U^2+U+I=0$, and hence
\[
 \Sigma_B^2=-3I,\qquad \Sigma_B^\dagger=I+2U^{-1}=I+2U^2=-\Sigma_B,\qquad \Sigma_B^\dagger \Sigma_B=3I.
\]

\begin{proposition}[Polarization-compatible $K$-action]\label{prop:Kcompatible}
The assignment $\sqrt{-3}\mapsto \Sigma_B$ defines an embedding
\[
 K=\Ql(\sqrt{-3})\hookrightarrow\operatorname{End}^0(A_z)
\]
at the certified $B_t$-vertex and is compatible with $E_1$: for every $k\in K$,
\[
 \eta(k)^\dagger\eta(k)=N_{K/\Ql}(k)I.
\]
\end{proposition}
\begin{proof}
The identity $\Sigma_B^2=-3I$ gives the field embedding. Since the vertex is fixed by $B_t$, the actual homology action is $B_t^{\mathsf T}$; every polynomial in it, in particular $\Sigma_B$, preserves the complex structure. Integrality gives an actual endomorphism. The relation $\Sigma_B^\dagger=-\Sigma_B$ gives
\[
 (a-b\Sigma_B)(a+b\Sigma_B)=a^2+3b^2=N_{K/\Ql}(a+b\sqrt{-3}),
\]
and the claim follows for all $a,b\in\Ql$.
\end{proof}


